\documentclass[aspectratio=169,11pt]{beamer}
\usepackage[T1]{fontenc}
\usepackage[utf8]{inputenc}
\usepackage{graphicx}
\usepackage{xcolor}
\usepackage{lmodern}

% Orange matching original slides
\definecolor{slideOrange}{RGB}{232,104,26}
\definecolor{hlYellow}{RGB}{255,235,130}

\usetheme{default}

\setbeamercolor{frametitle}{fg=slideOrange,bg=white}
\setbeamercolor{background canvas}{bg=white}
\setbeamercolor{normal text}{fg=black}
\setbeamercolor{itemize item}{fg=slideOrange}
\setbeamercolor{itemize subitem}{fg=slideOrange}
\setbeamercolor{enumerate item}{fg=slideOrange}
\setbeamercolor{title}{fg=slideOrange,bg=white}
\setbeamercolor{subtitle}{fg=darkgray,bg=white}

\setbeamerfont{frametitle}{size=\Large,series=\bfseries}
\setbeamerfont{title}{size=\huge,series=\bfseries}

\setbeamertemplate{navigation symbols}{}
\setbeamertemplate{itemize item}{\textbullet}
\setbeamertemplate{itemize subitem}{--}

\setbeamertemplate{footline}{%
  \vspace{2pt}%
  \noindent\textcolor{gray!40}{\rule{\paperwidth}{0.4pt}}%
  \vspace{3pt}%
}

\newcommand{\hly}[1]{\colorbox{hlYellow}{#1}}

\title{Introduction to Discrete Mathematics}
\subtitle{Lesson 1: What Is Discrete Mathematics?}
\author{Wolfram U}
\date{}

\begin{document}

% ── Title ────────────────────────────────────────────────────────────────────
{
\setbeamertemplate{footline}{}
\begin{frame}
  \centering
  \vspace{2.5em}
  {\usebeamerfont{title}\usebeamercolor[fg]{title}
    Introduction to Discrete Mathematics\par}
  \vspace{0.6em}
  {\large Lesson 1: What Is Discrete Mathematics?\par}
  \vspace{0.6em}
  {\normalsize\color{gray} Wolfram U\par}
\end{frame}
}

% ── The World of Integers ────────────────────────────────────────────────────
\begin{frame}{The World of Integers}
  \begin{columns}[T]
    \begin{column}{0.52\textwidth}
      \vspace{0.4em}
      Discrete mathematics is the study of \textbf{countable} things.

      \medskip
      It is not a branch of mathematics.

      \medskip
      Rather, it is a group of branches with a common property:
      the study of \textbf{discrete objects}.
    \end{column}
    \begin{column}{0.46\textwidth}
      \includegraphics[width=\textwidth]{slide_0006.png}
    \end{column}
  \end{columns}
\end{frame}

% ── Main Questions ───────────────────────────────────────────────────────────
\begin{frame}{Main Questions}
  Discrete mathematics aims to answer a variety of problems:

  \bigskip
  \begin{itemize}
    \setlength\itemsep{0.6em}
    \item What is \hly{known}?
    \item What can be \textit{computed}?
    \item What can be computed \textit{efficiently}?
    \item How can \textit{information} be organized and stored?
  \end{itemize}
\end{frame}

% ── Information ──────────────────────────────────────────────────────────────
\begin{frame}{Information}
  \begin{columns}[T]
    \begin{column}{0.54\textwidth}
      \vspace{0.4em}
      \hly{Logic, discrete structures and graphs} all allow the
      verification and storage of information.
    \end{column}
    \begin{column}{0.44\textwidth}
      \includegraphics[width=\textwidth]{slide_0008.png}
    \end{column}
  \end{columns}
\end{frame}

% ── Summarization ────────────────────────────────────────────────────────────
\begin{frame}{Summarization}
  \begin{columns}[T]
    \begin{column}{0.54\textwidth}
      \vspace{0.4em}
      Combinatorics, set theory and mathematical proofs allow the
      study of \textbf{emerging properties}.
    \end{column}
    \begin{column}{0.44\textwidth}
      \includegraphics[width=\textwidth]{slide_0010.png}
    \end{column}
  \end{columns}
\end{frame}

% ── Computation ──────────────────────────────────────────────────────────────
\begin{frame}{Computation}
  \begin{columns}[T]
    \begin{column}{0.54\textwidth}
      \vspace{0.4em}
      From a set of rules, \textbf{iterative and recursive
      computations} can be done and verified.
    \end{column}
    \begin{column}{0.44\textwidth}
      \includegraphics[width=\textwidth]{slide_0012.png}
    \end{column}
  \end{columns}
\end{frame}

% ── History ──────────────────────────────────────────────────────────────────
\begin{frame}{History}
  Integers were first used with \textbf{tally marks} (3400\,BC),
  and negative numbers were introduced with the
  \textbf{Chinese number rod system} (200\,BC).

  \medskip
  Many mathematicians have had an important influence on modern
  discrete mathematics, including:

  \bigskip
  \begin{columns}[c]
    \begin{column}{0.19\textwidth}
      \centering{\small\textbf{Aristotle}\\[2pt]Logic}
    \end{column}
    \begin{column}{0.19\textwidth}
      \centering{\small\textbf{L.\ Euler}\\[2pt]Graphs}
    \end{column}
    \begin{column}{0.19\textwidth}
      \centering{\small\textbf{G.\ Cantor}\\[2pt]Sets}
    \end{column}
    \begin{column}{0.19\textwidth}
      \centering{\small\textbf{P.\ Erd\H{o}s}}
    \end{column}
    \begin{column}{0.19\textwidth}
      \centering{\small\textbf{D.\ Knuth}\\[2pt]Algorithms}
    \end{column}
  \end{columns}
\end{frame}

% ── Discrete Mathematics Today ───────────────────────────────────────────────
\begin{frame}{Discrete Mathematics Today}
  Nowadays, discrete mathematics is strongly associated with
  \textbf{technology}.

  \bigskip
  \begin{itemize}
    \setlength\itemsep{0.5em}
    \item It is at the center of all \textbf{modern computation}.
    \item It is the cornerstone of
          \textbf{communications, networks and the internet}.
    \item Practically all types of \textbf{resource management}
          need discrete mathematics.
    \item Discrete math is important in
          \textbf{computer graphics, games} and even arts.
  \end{itemize}
\end{frame}

% ── Why Study Discrete Mathematics? ─────────────────────────────────────────
\begin{frame}{Why Study Discrete Mathematics?}
  High-school mathematics mostly overlooks discrete mathematics
  in favor of continuous mathematics.

  \medskip
  \textbf{Reasons to study discrete mathematics:}

  \begin{itemize}
    \setlength\itemsep{0.4em}
    \item Calculate discrete phenomena \textbf{efficiently}
    \item Understand \textbf{computation}
    \item \textbf{Use and store} information
    \item Develop your \textbf{mathematical reasoning}
  \end{itemize}
\end{frame}

% ── About This Course ────────────────────────────────────────────────────────
\begin{frame}{About This Course}
  This course is designed to be an introduction to many fields of
  mathematics and computer science.

  \medskip
  This exploration is facilitated by the use of
  \textbf{Wolfram Language}, a language with discrete computation
  as its core philosophy.

  \medskip
  This course is separated into \textbf{seven} small sections:

  \begin{columns}[T]
    \begin{column}{0.48\textwidth}
      \begin{enumerate}
        \item What Is Discrete Mathematics?
        \item Logic
        \item Sets
        \item Graphs
      \end{enumerate}
    \end{column}
    \begin{column}{0.48\textwidth}
      \begin{enumerate}
        \setcounter{enumi}{4}
        \item Combinatorics
        \item Number Theory
        \item Algorithms
      \end{enumerate}
    \end{column}
  \end{columns}
\end{frame}

% ── Who Is This Course For? ──────────────────────────────────────────────────
\begin{frame}{Who Is This Course For?}
  The Wolfram U course on discrete mathematics is most useful for:

  \bigskip
  \begin{itemize}
    \setlength\itemsep{0.5em}
    \item Future, current or self-taught
          \textbf{computer science and software engineering students}
    \item Individuals aiming to expand their
          \textbf{mathematical intuition} and problem-solving skills
    \item \textbf{Wolfram Language users} who want a guided
          introduction to discrete mathematics functions
    \item \textbf{Math hobbyists} wanting to explore the
          beauty of mathematics
  \end{itemize}
\end{frame}

% ── Summary ──────────────────────────────────────────────────────────────────
\begin{frame}{Summary}
  \begin{itemize}
    \setlength\itemsep{0.8em}
    \item As a \textbf{mathematical subject}, discrete mathematics
          is essentially the study of \textbf{countable objects}.

    \item As a \textbf{professional skill}, it helps you analyze,
          store and compute information, all to conquer many tasks.

    \item As an \textbf{educational hurdle}, it provides the
          toolset necessary to
          \hly{excel in computer science courses}.
  \end{itemize}
\end{frame}

\end{document}
